% chapters/08/section-06.tex
\section{AI Reaktif dan AI Berbasis Perencanaan}
\label{sec:reaktif-perencanaan}

NPC pada bab ini disebut reaktif karena seluruh keputusannya diambil dari
keadaan yang teramati saat itu. Bagian ini membahas batas kemampuan cara itu,
sesuai indikator \ind{c} pada bab ini.

\subsection{Cara Kerja AI Reaktif}

Kata ``reaktif'' berarti: keputusan pada bingkai ini hanya bergantung pada
keadaan pada bingkai ini. NPC tidak mengingat apa yang terjadi sebelumnya dan
tidak memperhitungkan apa yang akan terjadi kemudian. Ia mengamati, memutuskan,
dan bertindak, lalu mengulanginya.

Kesederhanaan itu membawa tiga keuntungan. Pertama, kodenya pendek dan mudah
diperiksa. Kedua, perhitungannya murah, sehingga jumlah NPC dapat banyak.
Ketiga, perilakunya mudah diramalkan oleh perancang, sehingga mudah diatur.

\subsection{Kelemahan yang Langsung Terasa}

Kelemahannya muncul begitu ada rintangan di antara NPC dan pemain. Perhatikan
apa yang diketahui NPC: hanya arah menuju pemain. Ia tidak mengetahui bahwa
ada dinding di antaranya, sebab dinding tidak termasuk dalam persepsinya.

Akibatnya NPC bergerak lurus ke arah pemain, menabrak dinding, dan tetap
menekan dinding itu selama pemain berada di baliknya. Dari sisi pemain, NPC
tampak bodoh: ia berhenti di depan dinding tanpa berusaha mencari jalan lain.
Keadaan itu diperlihatkan pada Gambar~\ref{fig:reaktif-perencanaan}.

\begin{figure}[htbp]
\centering
% chapters/08/figures/reaktif-vs-perencanaan.tikz
% Dua cara menyusun keputusan NPC: reaktif yang hanya melihat arah pemain,
% dan berbasis perencanaan yang memperhitungkan rintangan.
\begin{tikzpicture}[
  npc/.style={rectangle, draw=ubRed, thick, fill=ubRed!25, inner sep=4pt},
  pemain/.style={rectangle, draw=ubPrimary, thick, fill=ubPrimary!25, inner sep=4pt},
  dinding/.style={rectangle, draw=ubLine, thick, fill=ubGray},
  panah/.style={-{Stealth[length=2.2mm]}, very thick, draw=ubRed},
  panahjalan/.style={-{Stealth[length=2.2mm]}, very thick, draw=ubGreen},
  judul/.style={font=\footnotesize\bfseries, ubPrimary, align=center},
  label/.style={font=\scriptsize, align=center}
]

% ---------- Panel A: reaktif ----------
\begin{scope}
  \node[judul, anchor=south] at (2.4,2.35) {AI reaktif};
  \draw[dinding] (2.3,0.5) rectangle (2.65,2.3);

  \node[npc] (n1) at (0.8,1.4) {NPC};
  \node[pemain] (p1) at (4.0,1.4) {P};

  \draw[panah] (1.15,1.4) -- (2.2,1.4);
  \node[label, ubRed, anchor=north, text width=44mm] at (2.4,0.2)
    {bergerak lurus ke arah pemain, sehingga terus menabrak dinding dan tidak
     pernah sampai};
\end{scope}

% ---------- Panel B: berbasis perencanaan ----------
\begin{scope}[xshift=6.4cm]
  \node[judul, anchor=south] at (2.4,2.35) {AI berbasis perencanaan};
  \draw[dinding] (2.3,0.5) rectangle (2.65,2.3);

  \node[npc] (n2) at (0.8,1.4) {NPC};
  \node[pemain] (p2) at (4.0,1.4) {P};

  \draw[panahjalan] (1.15,1.4) -- (1.5,1.4);
  \draw[panahjalan] (1.5,1.4) -- (1.5,0.15) -- (3.5,0.15);
  \draw[panahjalan] (3.5,0.15) -- (3.5,1.4) -- (3.68,1.4);

  \node[label, ubGreen, anchor=north, text width=46mm] at (2.4,0.2)
    {memperhitungkan rintangan lebih dahulu, sehingga menemukan jalan
     memutar};
\end{scope}
\end{tikzpicture}

\caption{Perilaku reaktif dan berbasis perencanaan ketika ada rintangan}
\label{fig:reaktif-perencanaan}
\end{figure}

Satu hal yang perlu ditekankan: NPC itu tidak melakukan kesalahan. Ia
melaksanakan keputusannya dengan benar. Yang kurang adalah
\textbf{keterangan}---data yang tersedia baginya untuk mengambil keputusan.
Menambahkan keterangan itulah yang mengubah AI reaktif menjadi AI berbasis
perencanaan.

\subsection{Cara Kerja AI Berbasis Perencanaan}

AI berbasis perencanaan memperhitungkan lebih banyak keterangan sebelum
memutuskan. Pada bentuk yang paling sederhana, ia memeriksa apakah jalan
lurusnya terhalang, lalu mencari jalan memutar. Pada bentuk yang lebih umum,
ia menyusun rangkaian langkah menuju tujuan, menilai biayanya, dan memilih yang
paling murah.

Biaya perhitungannya jauh lebih besar, sebab mencari jalan memutar menuntut
penelusuran banyak kemungkinan. Karena itu, perencanaan biasanya tidak
dijalankan setiap bingkai, melainkan hanya ketika keadaan berubah, misalnya
ketika tujuan berpindah atau ketika jalan yang sedang ditempuh tidak lagi
berlaku. Cara menghemat perhitungan itu dibahas pada Bab~XI.

\subsection{Kapan Masing-Masing Dipakai}

Pilihan di antara keduanya bukan soal mana yang lebih baik, melainkan soal
kecukupan. Tabel~\ref{tab:pilihan-ai} memuat pertimbangannya.

\begin{table}[htbp]
\centering
\caption{Pertimbangan memilih cara pengambilan keputusan pada NPC}
\label{tab:pilihan-ai}
{\setlength{\tabcolsep}{4pt}\footnotesize%
\begin{tabularx}{\textwidth}{@{}L{4.4cm} L{3.0cm} Y@{}}
\toprule
\headrow \textbf{Keadaan permainan} & \textbf{Cara yang memadai} & \textbf{Alasan} \\
\midrule
Ruang terbuka tanpa rintangan &
Reaktif &
Tidak ada yang menghalangi, sehingga arah menuju pemain selalu dapat ditempuh.
Inilah keadaan pada bab ini. \\
\addlinespace
Ruang dengan rintangan kecil yang dapat dihindari dengan mudah &
Reaktif, ditambah penghindaran sederhana &
Perencanaan penuh terlalu mahal untuk rintangan yang jarang. \\
\addlinespace
Labirin atau ruang dengan banyak rintangan &
Berbasis perencanaan, yaitu pencarian jalur &
Jalan lurus hampir selalu terhalang, sehingga diperlukan pencarian rute.
Dibahas pada Bab~X. \\
\addlinespace
Banyak NPC yang harus berpindah tempat secara terkoordinasi &
Berbasis perencanaan, dengan penghematan perhitungan &
Biaya pencarian jalur menjadi pertimbangan utama. Dibahas pada Bab~XI. \\
\bottomrule
\end{tabularx}}
\end{table}

\subsection{Melengkapi Dua Kekurangan}

NPC pada bab ini memiliki dua kekurangan yang akan dilengkapi pada bab-bab
berikutnya.

\begin{itemize}
  \item \textbf{Kekurangan pengetahuan tentang rintangan.} Pada Bab~X, NPC
        akan memperoleh kemampuan mencari jalan memutar melalui pencarian jalur
        A*. Sebelum itu, Bab~IX melengkapinya dengan kemampuan menilai
        keadaan, sehingga NPC dapat memutuskan untuk berpatroli atau mengelak,
        bukan hanya mengejar atau berdiam.
  \item \textbf{Kekurangan pertimbangan biaya.} Pada bab ini, seluruh NPC
        menghitung persepsinya setiap bingkai. Selama jumlahnya sedikit, cara
        itu tidak menimbulkan masalah. Pada Bab~XI, ketika jumlah NPC dan
        luasnya pencarian jalur bertambah, biaya perhitungan menjadi
        pertimbangan yang menentukan.
\end{itemize}

\begin{catatabox}
Kedua kekurangan itu berbeda jenisnya. Yang pertama berkaitan
dengan \textbf{kualitas keputusan}, sedangkan yang kedua berkaitan dengan
\textbf{biaya keputusan}. Pada pengembangan permainan yang sesungguhnya, kedua
pertimbangan itu selalu berjalan bersama: keputusan yang lebih baik biasanya
menuntut perhitungan yang lebih mahal, dan tugas perancang adalah memilih
titik keseimbangan yang tepat.
\end{catatabox}
